\chapter{El producto operativo}
\label{cap:producto}

Un modelo que vive en un cuaderno de trabajo no le sirve a una cooperativa pesquera. Este
capítulo describe cómo se entrega el resultado del Capítulo~\ref{cap:resultados}: una
aplicación de inferencia que sirve, por especie y \ue{}, el intervalo calibrado de captura
diaria.

\section{Qué entrega y por qué así}

El producto \textbf{no es un total puntual por temporada}, y la decisión es deliberada. Como la
captura se concentra en pocos días alrededor de la apertura, la suma de las medianas diarias es
un mal estimador del total: subestima las series con picos grandes y sobreestima las de nivel
bajo. Lo que sí está calibrado ---y es lo que el sistema entrega--- es el \textbf{intervalo
diario} con su cobertura empírica verificada, más un resumen por temporada que muestra la
captura observada y la cobertura alcanzada, con la estimación central únicamente como
referencia. Es una diferencia de honestidad, no de presentación: se entrega la cantidad cuya
garantía se puede sostener.

\section{Arquitectura}

El sistema tiene tres piezas:

\begin{itemize}
  \item Una \textbf{API} (FastAPI) que, al arrancar, entrena una sola vez el modelo de
  producción ---el mismo de la Sección~\ref{sec:res_cqr}: \emph{gradient boosting} cuantílico
  en escala logarítmica con conformalización Mondrian por serie--- y deja en memoria, por serie,
  el pronóstico diario calibrado (mediana y bandas de 80 y 90\%). Los endpoints solo leen esa
  memoria, de modo que la consulta es inmediata. El entrenamiento inicial tarda entre 30 y 60
  segundos y el estado se expone en un endpoint de salud. Para despliegues reproducibles ese
  resultado puede además serializarse como \textbf{artefacto versionado}
  (\texttt{models/final/store.json}, unos 0.8 MB): es el pronóstico ya calculado ---no un
  \emph{pickle} del modelo, que ataría el despliegue a versiones exactas de las bibliotecas---
  acompañado de un manifiesto con el corte, la semilla, la versión del formato y la huella del
  conjunto de datos con el que se generó. Si el artefacto está presente la API arranca al
  instante y sirve exactamente las cifras reportadas aquí; si no, entrena en caliente y lo
  advierte en la bitácora.
  \item Un \textbf{front} estático mínimo: selector de especie y de zona, gráfico de abanico
  diario (observado contra las bandas calibradas) y un panel de \emph{backtest} por temporada
  con la captura observada y la cobertura empírica del intervalo del 90\% en temporada.
  \item Un \textbf{empaquetado} en Docker, de modo que el despliegue completo es un solo
  comando y no depende del entorno de desarrollo de quien lo corre.
\end{itemize}

El corte de prueba que define el horizonte servido es un parámetro de despliegue; por omisión se
usa 01-junio-2024, que es el corte con la calibración estable
(Sección~\ref{sec:res_moredata}). En su configuración actual sirve \textbf{28 series}: las 21 de
langosta roja más las de abulón y erizo.

\section{Lo que ve una cooperativa}

La Figura~\ref{fig:producto_api} muestra la salida de la API para tres series de langosta
elegidas por contraste, en la temporada 2024-2025.

\begin{figure}[htbp]
\centering
\includegraphics[width=\linewidth]{producto_api.png}
\caption{Intervalos calibrados servidos por la API para tres series de langosta roja en la
temporada 2024-2025: la serie insignia de San Quintín, una zona de historia larga (El Rosario) y
una de historia corta, incorporada solo con los avisos de CONAPESCA (Bahía Tortugas, Vizcaíno).
Los puntos son la captura observada; las bandas, los intervalos del 80 y 90\%; la línea, la
mediana. La cobertura empírica en temporada se anota en cada panel. El ancho de la banda
superior en los días de apertura es la limitación conocida del método: el cuantil alto en escala
logarítmica se amplifica al invertir la transformación.}
\label{fig:producto_api}
\end{figure}

Los tres paneles ilustran el rango de comportamiento del sistema. San Quintín, con el bache
dentro del entrenamiento, queda bien calibrada (92.5\%). El Rosario, con historia completa
2017-2026, alcanza 95.1\% y sus bandas siguen la estructura de picos de la temporada. Bahía
Tortugas, incorporada solo con avisos de CONAPESCA y con apenas dos temporadas de
entrenamiento, alcanza 93.2\% pero con bandas notoriamente más anchas: el sistema \emph{sabe}
que sabe menos de esa zona, y lo comunica ensanchando el intervalo en lugar de dar un número
falsamente preciso. Esa es exactamente la propiedad que se buscaba al reformular el pronóstico
en términos probabilísticos.

\section{Advertencias que el sistema declara}

Tres advertencias están documentadas y se muestran junto al resultado, porque un producto de
pronóstico que no declara sus límites induce confianza injustificada:

\begin{enumerate}
  \item La cobertura que reporta la interfaz es \textbf{en temporada}. Fuera de temporada la
  captura es cero y cualquier intervalo la cubre trivialmente, lo que inflaría el número.
  \item Las zonas de \textbf{historia corta} (Vizcaíno, Abreojos--La Purísima, incorporadas solo
  con datos de 2022 en adelante) están peor calibradas que las de historia larga. Su calibración
  depende de acumular temporadas, que es la conclusión de fondo de este trabajo.
  \item El sistema \textbf{no pronostica la oceanografía futura}. El horizonte llega hasta la
  última fecha con covariables disponibles: es una herramienta de pronóstico calibrado sobre el
  periodo con datos, no un oráculo de temporadas futuras arbitrarias.
\end{enumerate}
